% !TeX root = ../main.tex
\section{Glossary}
\label{app:glossary}

\begin{description}[style=unboxed,leftmargin=0cm,itemsep=3pt]
\item[Analysis by synthesis] Estimating settings by rendering candidates through
  the forward model and comparing them with the target.
\item[Cabinet IR] The impulse response of a speaker cabinet and microphone,
  applied by convolution.
\item[Capture] A neural model of one device at one setting, learned from
  dry/wet pairs (NAM, Kemper, TONEX).
\item[Cepstrum] Inverse Fourier transform of the log magnitude spectrum; used to
  build minimum-phase filters.
\item[CMA-ES] Covariance Matrix Adaptation Evolution Strategy, a gradient-free
  optimiser (\S\ref{sec:search}); \emph{separable}: diagonal covariance only.
\item[Convolution] Applying an impulse response to a signal.
\item[DDSP] Differentiable digital signal processing: DSP blocks written so that
  gradients flow through them.
\item[DI] Direct input: the dry guitar signal before any processing.
\item[Dry / wet] Unprocessed / processed signal.
\item[EQ allowance] The smooth, capped EQ the app's search distance forgives
  (\S\ref{sec:app-distance}).
\item[ESR] Error-to-signal ratio, the usual training loss of amp models.
\item[FAD] Fréchet Audio Distance between the embedding distributions of two
  sets of clips.
\item[Forward / inverse problem] Settings $\to$ sound / sound $\to$ settings.
\item[Grey-box model] A model with the circuit's structure and learned values.
\item[Impulse response (IR)] The output of an LTI system for a single click;
  describes the system completely.
\item[LFO] Low-frequency oscillator driving a modulation effect.
\item[Long-term average spectrum] A signal's spectrum averaged over a long
  stretch; differences between two of them give a matching EQ.
\item[LTI] Linear and time-invariant.
\item[Minimum phase] Of all filters with a given magnitude response, the one
  with the least delay.
\item[MR-STFT] Multi-resolution STFT distance.
\item[MUSHRA] A standard listening-test method with hidden reference and anchor.
\item[NAM] Neural Amp Modeler, the open capture format the app's NAM pedal plays.
\item[Pseudo-DI / stand-in dry signal] A dry signal not recorded with the target
  but made for it (here: from the notes).
\item[RMS] Root mean square, a measure of level.
\item[SI-SDR] Scale-invariant signal-to-distortion ratio, in dB.
\item[Stem] One instrument's audio separated from a mix.
\item[Template chain] The fixed pedal chain the app's search fits
  (Table~\ref{tab:template}).
\item[Tone stack] The bass/middle/treble filter of an amplifier.
\item[Unpaired / wet only] Dry audio from another performance is available /
  no dry audio at all.
\end{description}
